\section{Appendix B: Spectral budget and assignment bounds}

The attaining assignment in \cref{obj:def:capacity-handle} combines a componentwise regularity bound with ordered spectral balance. The exact reflection inequality preserves the pooled budget when the nonperiodic prognostic function is represented on the torus. The common shift then converts expected prognostic imbalance into a sum of frequency-specific second moments. Finally, the rounding guarantee controls those moments while preserving fair assignment on the original units.

\subsection{Exact reflection and the pooled budget}

For the canonical components in \cref{obj:synth_1}, introduce the appendix-local reflected functions
\[
h_j(y)=g_j(m)(b_1(y)),
\qquad
h_{j\ell}(y)=g_{j\ell}(m)(b_2(y)),
\]
on \(\mathbb T_2^1\) and \(\mathbb T_2^2\), respectively. Their component Fourier coefficients, with normalized Haar integration, are
\[
\widehat h_j(a)
=\int_{\mathbb T_2^1}h_j(y)e^{-\mathrm i\pi ay}\,dy,
\qquad a\in\mathbb Z,
\]
and
\[
\widehat h_{j\ell}(a)
=\int_{\mathbb T_2^2}h_{j\ell}(y)e^{-\mathrm i\pi a^{\mathsf T}y}\,dy,
\qquad a\in\mathbb Z^2.
\]
These coefficients connect each component's restriction norm to its contribution to the full reflected spectrum. The following statement records the exact normalization before the broader transfer result.

\begin{lemmav}{lem:exact-reflection-budget}[Exact reflection budget]
Let \(p\), \(s\), and the real-valued cube function \(g\) satisfy:
\begin{itemize}
\item \textbf{(Dimension.)} \(p\in\{1,2\}\).
\item \textbf{(Smoothness.)} \(0<s\leq1\).
\item \textbf{(Restriction norm.)} \(g\in L^2([0,1]^p)\), with a Borel representative, and \(N_{p,s}(g)^2<\infty\), where
\[
N_{p,s}(g)^2
=
\inf_{\substack{G\in L^1(\mathbb R^p;\mathbb C)\cap L^2(\mathbb R^p;\mathbb C)\\
G=g\ \text{a.e. on }[0,1]^p}}
\int_{\mathbb R^p}
\left(1+\frac{\|\omega\|_2^2}{p}\right)^s
\left|
(2\pi)^{-p/2}
\int_{\mathbb R^p}G(u)e^{-\mathrm i\omega^{\mathsf T}u}\,du
\right|^2\,d\omega .
\]
\end{itemize}
On the period-two torus \(\mathbb T_2^p=(\mathbb R/2\mathbb Z)^p\), define the coordinatewise reflection map and reflected function by
\[
(b_p(y))_j=\min\{y_j,2-y_j\},
\qquad y\in[0,2)^p,
\qquad h(y)=g(b_p(y)).
\]
For \(a\in\mathbb Z^p\), define the normalized Fourier coefficient
\[
\widehat h(a)
=
2^{-p}\int_{[0,2)^p}h(y)e^{-\mathrm i\pi a^{\mathsf T}y}\,dy .
\]
Then
\[
\sum_{a\in\mathbb Z^p}
\left(1+\frac{\pi^2\|a\|_2^2}{p}\right)^s
|\widehat h(a)|^2
\leq N_{p,s}(g)^2 .
\]

For every integer \(d\geq2\), every \(0<s\leq1\), and every \(m\in\mathcal H_{d,s}\) as defined in \cref{obj:def:sobolev-class}, let \(F\) be its reflected function on \(\mathbb T_2^d\), with
\[
F(y)=m(b_d(y)),
\qquad
(b_d(y))_j=\min\{y_j,2-y_j\},
\qquad y\in[0,2)^d,
\]
and define
\[
\widehat F(k)
=
2^{-d}\int_{[0,2)^d}F(y)e^{-\mathrm i\pi k^{\mathsf T}y}\,dy,
\qquad
\operatorname{supp}(k)=\{j:k_j\neq0\},
\qquad k\in\mathbb Z^d.
\]
For nonzero \(k\), put
\[
t(k)=\frac{\|k\|_2^2}{|\operatorname{supp}(k)|}.
\]
Then
\[
\widehat F(0)=0,
\qquad
\widehat F(k)=0
\quad\text{whenever }|\operatorname{supp}(k)|>2,
\]
and
\[
\sum_{\substack{k\in\mathbb Z^d\\1\leq|\operatorname{supp}(k)|\leq2}}
(1+\pi^2t(k))^s|\widehat F(k)|^2\leq1,
\qquad
\int_{[0,1]^d}|m(x)|^2\,dx\leq1.
\]
\end{lemmav}

The component inequality in \cref{obj:lem:exact-reflection-budget} preserves the squared restriction budget with constant one throughout \(0<s\leq1\). Canonical centering assigns main effects to frequencies supported on their single coordinate and pair effects to frequencies supported on both of their coordinates. Thus the full spectral budget retains the component allocation specified by \cref{obj:ass:pooled-budget}. Its exact normalization supplies the budget used in the numerical attaining bound.

The accompanying support and centering conclusions identify the spectral rows relevant to the population criterion. The reflected function has frequencies supported on one or two coordinates and has zero constant coefficient. Its cube \(L^2\) bound also controls the independent-sign branch of \cref{obj:def:capacity-handle}.

\subsection{Common-shift spectral transfer}

The randomized lift in \cref{obj:synth_2} makes each lifted covariate Haar-uniform. A common translation preserves the product Haar law and makes the transformed sample independent of the shift. This independence is the basis of the spectral identity: the shift supplies a common random phase for each frequency while assignment depends on the transformed sample. The next statement joins the reflection bound to this identity for the allowed measurable prognostic functions.

\begin{lemmav}{lem:reflection-transfer}[Reflection and spectral transfer]
Write \(\mathbb T_2^p=(\mathbb R/2\mathbb Z)^p\) for the period-two torus, equipped with normalized Haar measure, and let \(b_p\) fold each coordinate onto \([0,1]\), by taking its absolute value in a representative from \([-1,1]\). For a cube function \(g\), its even reflection \(h\) and normalized Fourier coefficients are
\[
h=g\circ b_p,\qquad
\widehat h(a)=\int_{\mathbb T_2^p}h(w)
  \exp(-\mathrm i\pi a\cdot w)\,dw,
\qquad a\in\mathbb Z^p.
\]
The squared restriction norm is
\[
N_{p,s}(g)^2
=
\inf_{\substack{G\in L^1(\mathbb R^p;\mathbb C)\cap L^2(\mathbb R^p;\mathbb C)\\
G=g\ {\rm a.e.}\ {\rm on}\ [0,1]^p}}
\int_{\mathbb R^p}
\left|
(2\pi)^{-p/2}
\int_{\mathbb R^p}\exp(-\mathrm i\omega\cdot u)G(u)\,du
\right|^2
\left(1+\frac{\|\omega\|^2}{p}\right)^s\,d\omega.
\]
For every \(p\in\{1,2\}\) and \(0<s\leq1\), there exists a constant \(c>0\), depending on \(p,s\), such that every real cube function \(g\), under the usual measurability and square-integrability conditions and with \(N_{p,s}(g)^2<\infty\), satisfies
\[
\sum_{a\in\mathbb Z^p}
\left(1+\frac{\pi^2\|a\|^2}{p}\right)^s
|\widehat h(a)|^2
\leq c\,N_{p,s}(g)^2.
\]

For every \(d\geq2\), \(0<s\leq1\), and \(m\in\mathcal H_{d,s}\) as in \cref{obj:def:sobolev-class}, define the reflected function \(F\), its Fourier coefficients, and the frequency complexity by
\[
F=m\circ b_d,\qquad
\widehat F(k)=\int_{\mathbb T_2^d}F(w)
  \exp(-\mathrm i\pi k\cdot w)\,dw,
\]
\[
\operatorname{supp}(k)=\{j:k_j\neq0\},\qquad
t(k)=
\begin{cases}
\displaystyle\frac{\sum_j k_j^2}{|\operatorname{supp}(k)|},&k\neq0,\\
0,&k=0.
\end{cases}
\]
Then, with \(C_s=3072\),
\[
\sum_{\substack{k\in\mathbb Z^d\\
1\leq|\operatorname{supp}(k)|\leq2}}
(1+\pi^2t(k))^s|\widehat F(k)|^2
\leq C_s.
\]

For every pair of nonnegative integers \(n,d\), let \(X\) satisfy \cref{obj:ass:uniform-draw} on an arbitrary measurable sampling space. Independently of \(X\), reflect each coordinate of each \(X_i\) by an independent fair sign, and add an independent common Haar shift \(T\in\mathbb T_2^d\), obtaining \(W_i\) modulo \(2\). The joint law of \(W=(W_i)_{i=1}^n\) and \(T\) is the product of \(n+1\) normalized Haar laws. Thus the \(W_i\) are independent Haar draws and \(W\) is independent of \(T\).

The following identity holds under these conditions:
\begin{itemize}
\item \textbf{(Parameters.)} \(n\geq2\), \(d\geq2\), and \(0<s\leq1\).
\item \textbf{(Prognostic function.)} \(m\in\mathcal H_{d,s}\) as in \cref{obj:def:sobolev-class}.
\item \textbf{(Sampling.)} \(X\) satisfies \cref{obj:ass:uniform-draw} on an arbitrary measurable sampling space, and \(W,T\) are constructed as above.
\item \textbf{(Assignment.)} \(\pi\) is any Borel probability kernel from \((\mathbb T_2^d)^n\) to \(\mathcal Z_n=\{-1,1\}^n\), and, conditional on the sample, reflection signs, and shift, \(Z\) has law \(\pi(W)\).
\end{itemize}
With the torus character \(e_k(w)=\exp(\mathrm i\pi k\cdot w)\), one has
\[
\mathbb E\left|\sum_{i=1}^n Z_i m(X_i)\right|^2
=
\sum_{k\in\mathbb Z^d}
|\widehat F(k)|^2\,
\mathbb E\left|\sum_{i=1}^n Z_i e_k(W_i)\right|^2.
\]
On the left, expectation includes the original sample, independent reflection signs, common shift, and conditional assignment law; on the right, it includes the independent Haar sample \(W\) and conditional assignment law \(\pi(W)\). The equality is understood in the extended nonnegative reals.
\end{lemmav}

The transfer in \cref{obj:lem:reflection-transfer} separates population regularity from assignment performance. Squared Fourier coefficients describe the reflected prognosis; the expected squared signed character sums describe the assignment rule. Orthogonality of the common-shift phases removes cross-frequency terms, even when the assignment signs depend on the entire transformed sample. The dependence condition is therefore substantive: conditional assignment uses \(W\), while the independent shift supplies the averaging that diagonalizes the criterion.

The identity applies to every Borel representative allowed by \cref{obj:def:sobolev-class}. Its square-integrable formulation accommodates Fourier expansions in the \(L^2\) sense, and uniform sampling makes changes on cube null sets immaterial to the original-sample loss. Reflection consequently carries nonperiodic component regularity into the spectral assignment criterion with the boundary behavior permitted by the restriction norm. For the attaining comparison, the sharper unit budget in \cref{obj:lem:exact-reflection-budget} supplies the normalization within this transfer.

\subsection{Ordered rounding and the attaining bound}

The spectral rows are bounded real functions ordered by increasing frequency complexity. Their assignment guarantees must hold conditional on each realized input matrix, so that auxiliary transformations can subsequently be averaged. The following statement gives the measurable realization, finite completion, fairness, and second-moment bounds for the procedure in \cref{obj:def:ordered-boundary-rounding}.

\begin{lemmav}{lem:ordered-prefix-rounding}[Ordered prefix rounding guarantees]
Let \(n\geq0\) be an integer, let \(\Lambda\in\mathbb R\), and let \(a_1,a_2,\ldots\in\mathbb R^n\) be deterministic ordered rows satisfying
\[
|(a_h)_i|\leq\Lambda
\qquad(h\geq1,\ 1\leq i\leq n).
\]
Supply the first \(K=\lfloor n/4\rfloor\) rows to the procedure in \cref{obj:def:ordered-boundary-rounding}, and write \(Z\in\{-1,1\}^n\) for its output signs. The output map is jointly Borel measurable in the input matrix \(a\in\mathbb R^{K\times n}\) and the real seed vector of length \(3n\). For every seed vector in \([0,1]^{3n}\), the fractional state after \(n\) moves satisfies
\[
u_i=Z_i\qquad(1\leq i\leq n).
\]
With the \(3n\) seeds drawn independently and uniformly from \([0,1]\),
\[
\mathbb E[Z_i]=0
\qquad(1\leq i\leq n),
\qquad
\mathbb E\!\left[(a_h^{\mathsf T}Z)^2\right]
\leq32\Lambda^2\min\{h,n\}
\qquad(h\geq1).
\]
\end{lemmav}

Joint Borel measurability in \cref{obj:lem:ordered-prefix-rounding} makes the finite randomized map suitable for covariate-dependent assignment. The ordered projection and basis conventions in \cref{obj:def:ordered-boundary-rounding} specify its choices on matrices of every rank, including boundary cases. Completion holds for every allowed seed vector; fairness and the second-moment bounds concern independent uniform seeds. Integrating this map over the rounding seeds therefore gives a Borel assignment law with fair marginals for every fixed input matrix.

The boundary probabilities preserve coordinate means, while the active constraints protect an ordered prefix of rows. Earlier rows receive longer protection as the active set contracts. The bound \(32\Lambda^2\min\{h,n\}\) quantifies this allocation and applies to the entire ordered row family, including rows beyond the supplied prefix. This last feature connects the finite input matrix to the full spectral budget.

For the prescribed cosine and sine rows, the entry bound is \(\Lambda=\sqrt2\). Conditional application of \cref{obj:lem:ordered-prefix-rounding} controls their signed second moments, and \cref{obj:lem:reflection-transfer} transfers those controls to the original prognostic imbalance. Ordering by frequency complexity aligns the strongest row bounds with the spectral weights in \cref{obj:lem:exact-reflection-budget}. Together, these are the assignment and regularity ingredients of the spectral-branch comparison in \cref{obj:thm:log-free-bivariate-capacity}.

The construction remains admissible after averaging over independent reflection coins and the common shift. Conditional fairness holds for every transformed input, and the randomization uses the original covariates and independent auxiliary seeds as specified in \cref{obj:def:capacity-handle}. Hence the resulting original-sample assignment belongs to \cref{obj:def:design-class}.

When \(n\leq B\), where \(B=\binom d2\) is the coordinate-pair count, the independent-sign branch is controlled by the cube \(L^2\) bound in \cref{obj:lem:exact-reflection-budget}. When \(B<n\), the ordered spectral branch gives the smoothness-dependent comparison. The assembled guarantee, stated in \cref{obj:thm:log-free-bivariate-capacity}, is
\[
\sup_{m\in\mathcal H_{d,s}}
\mathcal L_{n,d,s}(\pi^\star_{n,d},m)
\leq
3072\min\{1,(B/n)^s\},
\qquad n,d\geq2,\quad 0<s\leq1.
\]
Here the assignment \(\pi^\star_{n,d}\) is the construction in \cref{obj:def:capacity-handle}, and the loss is the criterion specified in \cref{obj:thm:log-free-bivariate-capacity}. Smoothness enters through the spectral weights used to evaluate this single assignment procedure, yielding simultaneous attainment across the stated smoothness range.
